\section{Einleitung}
\label{sec:einleitung}

Die fortschreitende Entwicklung autonomer Systeme prägt zunehmend technologische Innovationen in Industrie, Logistik und Mobilität. Insbesondere im Bereich der autonomen Fahrzeuge ermöglichen sie eine Effizienzsteigerung, erhöhte Sicherheit und die Erschließung neuer Anwendungsszenarien. Die Realisierung solcher Systeme auf ressourcenbeschränkter Hardware, wie sie in Modellfahrzeugen oder spezialisierten Robotern zum Einsatz kommt, stellt jedoch eine besondere Herausforderung dar. Sie erfordert eine sorgfältige Abstimmung von Sensorik, Aktorik, Softwarearchitektur und Regelungsalgorithmen, um echtzeitnahe und robuste Funktionalität zu gewährleisten.

Diese Studienarbeit widmet sich der Konzeption und Implementierung eines teilautonomen Fahrsystems auf Basis eines existierenden Modellfahrzeugs. Im Zentrum steht die Entwicklung einer lokalen Hindernisvermeidung unter Verwendung von \gls{lidar}-Sensorik und dem \gls{ros}. Die wissenschaftliche Problemstellung ergibt sich aus der Notwendigkeit, komplexe Wahrnehmungs- und Planungsalgorithmen so zu adaptieren, dass sie auf einer Embedded-Plattform wie dem Raspberry Pi performant ausgeführt werden können. Ein zentraler Untersuchungsgegenstand ist dabei die Eignung eines reaktiven, auf Gaußfunktionen basierenden Ansatzes zur Trajektorienplanung, der ohne globale Kartierung eine flüssige und sichere Umfahrung statischer Hindernisse ermöglichen soll.

Das Hauptziel dieser Arbeit ist die Weiterentwicklung des Modellfahrzeugs zu einem robusten, teilautonomen System, das sowohl manuell via Smartphone-Applikation als auch autonom betrieben werden kann. Die Funktionalität wird anhand messbarer Kriterien evaluiert, darunter die zuverlässige Erkennung von Hindernissen, die Erfolgsquote von Ausweichmanövern bei definierter Geschwindigkeit und die Schaffung einer reproduzierbaren Softwareumgebung mittels Containerisierung. Die Arbeit leistet damit einen Beitrag zur praktischen Anwendung und Bewertung von Algorithmen für das autonome Fahren im Kontext ressourcenlimitierter Systeme.